\section{Результаты}

Программа собрана и запущена под Linux (WSL2, Ubuntu). Замеры выполнялись на матрице
$2048\times2048$, окно 3, $K=10$ применений фильтра: \texttt{./lab3 --random --rows 2048
--cols 2048 --iters 10 --bench 16 --repeat 5 --check}. Для каждого $P\in\{1,2,4,8,16\}$
выполнялось 5 повторов, в таблице --- лучшее время; $S(P)=T_1/T_P$, $E(P)=100\cdot S(P)/P$.

\begin{center}
\begin{tabular}{rrrrrr}
\hline
$P$ & блоков & $T_P$ best, мс & $T_P$ med, мс & $S(P)$ & $E(P)$, \% \\
\hline
1  & 4  & 9013.429 & 9457.031 & 1.000 & 100.00 \\
2  & 8  & 6739.713 & 8158.216 & 1.337 & 66.87 \\
4  & 16 & 3551.429 & 3770.643 & 2.538 & 63.45 \\
8  & 32 & 1828.926 & 2064.036 & 4.928 & 61.60 \\
16 & 64 & 1095.323 & 1110.586 & 8.229 & 51.43 \\
\hline
\end{tabular}
\end{center}

Проверка параллельного результата последовательным эталоном во всех прогонах:
\texttt{check: OK}.

\textbf{Объяснение результатов.} До $P=8$ ускорение растёт почти линейно ($S(8)\approx 4.9$):
блоки не пересекаются, число блоков кратно числу потоков, и на синхронизацию уходит малая доля
времени --- критическая секция содержит только захват номера блока. Эффективность
$E\approx 63\%$ ниже $100\%$ из-за:
\begin{itemize}
  \item барьера на каждой итерации --- самый быстрый поток ждёт самый медленный; балансировка
    блоками по строкам не идеальна (краевой блок короче);
  \item алгоритм память-зависимый: окно $3\times3$ (сортировка 9 элементов) даёт низкое
    отношение вычислений к работе с памятью, потоки конкурируют за пропускную способность
    памяти и кэш;
  \item накладных расходов на создание потоков и удержание мьютекса.
\end{itemize}
При переходе от 8 к 16 потокам прирост замедляется ($S(16)\approx 8.2$, $E\approx 51\%$):
исполняемых аппаратно контекстов меньше 16, поэтому часть 16 потоков фактически делит
физические ядра, а издержки синхронизации (барьер с 16 участниками) растут.

\textbf{Демонстрация числа потоков штатными средствами ОС.} Запуск \texttt{--threads 6} на
большой матрице ($4096\times4096$, $K=20$): \texttt{ps -eLf} показывает 6 системных потоков
(LWP) одного процесса, \texttt{/proc/<pid>/status} --- \texttt{Threads: 6}, \texttt{ps -L} ---
\texttt{NLWP 6}: главный поток + 5 созданных \texttt{pthread\_create}, то есть лимит ключа
запуска соблюдается.

\begin{lstlisting}[caption={ps -eLf / /proc/<pid>/status при --threads 6},captionpos=b]
$ ./lab3 --random --rows 4096 --cols 4096 --iters 20 --threads 6 &
$ ps -eLf | grep "[l]ab3"
    PID     LWP C NLWP COMMAND
    632     632 92   6 ./lab3 --random --rows 4096 --cols 4096 --iters 20 --threads 6
    632     634 88   6 ./lab3 --random --rows 4096 --cols 4096 --iters 20 --threads 6
    632     635 88   6 ./lab3 --random --rows 4096 --cols 4096 --iters 20 --threads 6
    632     636 88   6 ./lab3 --random --rows 4096 --cols 4096 --iters 20 --threads 6
    632     637 88   6 ./lab3 --random --rows 4096 --cols 4096 --iters 20 --threads 6
    632     638 88   6 ./lab3 --random --rows 4096 --cols 4096 --iters 20 --threads 6
$ grep -i Threads /proc/632/status
Threads:  6
$ ps -o pid,lwp,nlwp,comm -L -p 632
    PID     LWP NLWP COMMAND
    632     632    6 lab3
    632     634    6 lab3
    632     635    6 lab3
    632     636    6 lab3
    632     637    6 lab3
    632     638    6 lab3
\end{lstlisting}
\pagebreak
